In order to balance computation time against generalization across diverse manipulation tasks, I developed an multi-query kinodynamic planner employing lazy evaluation techniques for collision-checking and dynamics propagation \cite{pasricha2024virtues}. Despite the significant improvement in computation times and success rates for planning time-parameterized trajectories, the approach remains inherently limited in scalability, especially when considering the combinatorial complexity of robot-environment interactions which is only compounded by inaccurate dynamics models.
Consequently, we conducted a broader investigation into transformer architectures for robot dynamics modeling, enabling accurate short-horizon state predictions across diverse robotic platforms \cite{reed2024transformer}. Through comparative studies on both manipulators and ground vehicles, we demonstrate superior generalization capabilities and enhanced interpretability through attention mechanisms compared to prior analytical and learned dynamics modeling approaches.
Additionally, we addressed the limitations imposed by overly conservative simplified dynamics models in the context of spill-free liquid handling (such as pendulum or mass-spring-damper approximations) \cite{muchacho2022solution} by introducing geometric inertial constraints to bias sampling during planning and a transformer-based model to assess whether a geometric path is likely to lead to a valid time-optimal time-parameterization before applying the parameterization itself \cite{abderezaei2024clutter}. Specifically demonstrated in spill-free liquid transport tasks, this strategy dramatically expanded the feasible solution space and significantly improved computational efficiency, reinforcing the utility of task-specific, data-driven dynamics modeling \cite{reed2024transformer}.